\section{整闭 Noether 整环上的模}

在本段中，A 始终表示一个分式域为 K 的交换整环。从第 2 节起，进一步假设 A 是 Noether 且整闭的（因而是 Krull 整环（§ 1，第 3 节，定理 2 的推论））；于是 P(A)、D(A) 和 C(A) 分别表示 A 的高度为 1 的素理想集合（§ 1，第 6 节）、A 的除子群（§ 1，第 3 节）和 A 的除子类群（§ 1，第 10 节），后两个群均用加法记号书写。

研究整闭 Noether 整环 A 上有限生成模的一般方法，是相对于 A 中所有高度为 1 的素理想 \(p \in \mathrm{P}(\mathrm{A})\) 对这些模作``局部化''；因为此时 \(A_{p}\) 是离散赋值环（§ 1，第 6 节，定理 4），有限生成 \(A_{p}\)-模的结构是众所周知的（代数，第七章，§ 4），因而能提供关于有限生成 A-模结构的信息。特别地，当 A 是 Dedekind 整环时，我们可以得到与 A 为主理想整环时（第 10 节）同样完备的理论。

